\documentclass[11pt]{article}
\usepackage[margin=1in]{geometry}
\usepackage[T1]{fontenc}
\usepackage{lmodern}
\usepackage{microtype}
\usepackage{amsmath,amssymb,amsthm,mathtools}
\usepackage{booktabs}
\usepackage{enumitem}
\usepackage{xcolor}
\usepackage[hidelinks]{hyperref}
\usepackage{fancyhdr}

\definecolor{notegray}{RGB}{245,246,248}
\newtheorem{theorem}{Theorem}
\newtheorem{proposition}[theorem]{Proposition}
\newtheorem{lemma}[theorem]{Lemma}
\newcommand{\permanent}{\operatorname{per}}
\newcommand{\R}{\mathbb{R}}

\setlength{\headheight}{14pt}
\pagestyle{fancy}
\fancyhf{}
\lhead{Dittert conjecture for $n=5$}
\rhead{Exact computer-assisted proof note}
\cfoot{\thepage}

\title{An Exact Computer-Assisted Proof of the\\Dittert Conjecture for $n=5$}
\author{Pedro Paulo Marques do Nascimento}
\date{23 July 2026\\[0.4em]
\small Licensed under \href{https://creativecommons.org/licenses/by/4.0/}{CC BY 4.0}}

\hypersetup{
  pdftitle={An Exact Computer-Assisted Proof of the Dittert Conjecture for n=5},
  pdfauthor={Pedro Paulo Marques do Nascimento},
  pdfsubject={Exact working proof and reproducible certificate},
  pdfkeywords={Dittert conjecture, permanent, computer-assisted proof, exact certificate}
}

\begin{document}
\maketitle

\begin{center}
\fcolorbox{black!25}{notegray}{\parbox{0.90\textwidth}{\small
\textbf{Review status.} This document reports a new exact computational
certificate produced in an interactive research session.  The complete
integer certificate and verifier are included.  The argument has not yet
undergone independent peer review and should be checked before being cited as
an established result.
}}
\end{center}

\begin{abstract}
For a nonnegative $5\times5$ matrix $A$ whose entries sum to $5$, let
\[
 \Phi(A)=\prod_{i=1}^{5}r_i+\prod_{j=1}^{5}c_j-\permanent(A),
\]
where $r_i$ and $c_j$ are the row and column sums.  We give an exact
computer-assisted proof that
\[
 \Phi(A)\le 2-\frac{5!}{5^5}=\frac{6130}{3125},
\]
with equality only at the matrix $J_5$ whose entries are $1/5$.  Published
structural results reduce the boundary case to the face with two zero entries
in distinct rows and columns.  On that face an integer
multiplier--sum-of-squares certificate proves a strict quintic inequality.
The verifier checks all $80\,730$ quintic coefficients using arbitrary-precision
integer arithmetic.  The smallest residual coefficient is approximately
$5.4206471\cdot10^{-5}$.
\end{abstract}

\section{Statement and boundary reduction}
Let
\[
 K_5=\left\{A=(a_{ij})\in\R_{\ge0}^{5\times5}:
             \sum_{i,j}a_{ij}=5\right\}.
\]

\begin{theorem}\label{thm:main}
For every $A\in K_5$,
\[
 \Phi(A)\le\frac{6130}{3125}.
\]
Equality holds if and only if $A=J_5$.
\end{theorem}

We use two published structural facts.  Hwang proved that the only possible
positive $\Phi$-maximizer is $J_n$ \cite{Hwang1986}.  Cheon and Wanless proved
that a partly decomposable matrix cannot maximize $\Phi$, and that the
complete zero set of a maximizer cannot, up to row and column permutations,
be one proper rectangular block \cite[Theorems~3.2 and~3.7]{CheonWanless2012}.

\begin{lemma}\label{lem:zeros}
If a nonempty set of matrix positions contains no two positions in distinct
rows and distinct columns, then all its positions lie in one row or all lie
in one column.
\end{lemma}

\begin{proof}
Choose $(i,j)$ in the set.  Every other position shares row $i$ or column $j$
with it.  A position of each type would give two positions in distinct rows
and columns, so only one type can occur.
\end{proof}

\begin{proposition}\label{prop:reduction}
It is enough to prove
\[
 \Phi(A)<\frac{6130}{3125}
 \qquad(A\in K_5,\ a_{11}=a_{22}=0).
 \tag{1}\label{eq:facegoal}
\]
\end{proposition}

\begin{proof}
Compactness gives a global maximizer.  A positive maximizer is $J_5$ by
Hwang's theorem.  If a boundary maximizer had no two independent zeroes,
Lemma~\ref{lem:zeros} would put its complete zero set in one row or column.
If the whole row or column vanished the matrix would be partly decomposable;
otherwise its zero set would be one proper rectangular block.  Both cases are
excluded by Cheon--Wanless.  Row and column permutations therefore place two
zeroes at $(1,1)$ and $(2,2)$, where \eqref{eq:facegoal} excludes the boundary
maximizer.
\end{proof}

\section{The quintic on the two-zero face}
Put $X=A/5$.  Its entries sum to $1$.  On the face $x_{11}=x_{22}=0$, index
the remaining $23$ entries by
\[
 V=\{(i,j):1\le i,j\le5\}\setminus\{(1,1),(2,2)\}.
\]
Write $S=\sum_{v\in V}x_v$.  Let $\mathcal E$ be the collection of
five-element subsets of $V$ that meet all five rows or all five columns, and
put
\[
 F(x)=\sum_{E\in\mathcal E}\prod_{v\in E}x_v.
 \tag{2}
\]
There are $2000$ row transversals, $2000$ column transversals, and $78$
allowed permutation transversals, hence $|\mathcal E|=3922$.

Expanding the row-sum and column-sum products shows that their common
square-free monomials are exactly the permanent terms.  Consequently
\[
 \prod_i R_i+\prod_j C_j-\permanent(X)=F(x),
 \qquad \Phi(A)=5^5F(x).
 \tag{3}
\]
Thus \eqref{eq:facegoal} is equivalent to
\[
 G(x):=\frac{1226}{5^9}S^5-F(x)>0
 \qquad(x\ge0,\ S=1).
 \tag{4}\label{eq:G}
\]

\section{Exact certificate}
Let $m_2(x)$ be the vector of all $276$ quadratic monomials in the $23$
variables.  The automorphism group $\Gamma$ of the two-zero face has order
$144$.  Its action has three vertex orbits, of sizes $2$, $12$, and $9$;
choose representatives $u_1,u_2,u_3$.

The certificate archive contains integer lower-triangular matrices
\[
 \widehat L_p\in\mathbb Z^{276\times276}\qquad(1\le p\le3),
\]
and sets $L_p=2^{-28}\widehat L_p$.  The verifier expands the identity
\begin{align}
 G(x)
 &=\sum_{p=1}^3\sum_{\gamma\in\Gamma}
 x_{\gamma u_p}
 \left\|L_p^{\mathsf T}m_2(\gamma x)\right\|_2^2
 +\frac1D\sum_{|\alpha|=5}R_\alpha x^\alpha,
 \tag{5}\label{eq:certificate}\\
 D&=5^9 2^{56}=140\,737\,488\,355\,328\,000\,000\,000.\notag
\end{align}
It checks all
\[
 \binom{23+5-1}{5}=80\,730
\]
coefficients and finds
\[
 \min_{|\alpha|=5}R_\alpha
 =7\,628\,882\,599\,067\,611\,080>0.
 \tag{6}
\]

\begin{proposition}\label{prop:face}
The strict inequality \eqref{eq:G} holds on the nonnegative simplex $S=1$.
\end{proposition}

\begin{proof}
Every term in the first sum of \eqref{eq:certificate} is nonnegative for
$x\ge0$: it is a nonnegative coordinate times a square norm.  The residual
polynomial has strictly positive coefficients by \eqref{eq:certificate}--(6).
Since $S=1$, at least one coordinate is positive, and its fifth-power term
makes the residual polynomial strictly positive.
\end{proof}

Proposition~\ref{prop:face} and Proposition~\ref{prop:reduction} prove
Theorem~\ref{thm:main}.

\section{Exact verification and scope}
The numerical semidefinite program was used only to discover the three Gram
matrices.  Each was shifted into the positive-definite cone, Cholesky factored,
and rounded on a $2^{-28}$ grid.  The saved integer factors are the proof
object; no floating-point number is used by the verifier for a correctness
decision.

The verifier independently reconstructs the face, the group, all $3922$
hyperedges, the three Gram matrices, and every coefficient of
\eqref{eq:certificate}.  Run, without Python's \texttt{-O} option,
\begin{verbatim}
python3 verify_primary.py dittert_n5_exact_certificate.npz
\end{verbatim}
The expected final line is \texttt{CERTIFIED}.  The certificate SHA-256 digest
is
\begin{center}
\texttt{373041bda29e1164059a2adbad1aaacd2911784a273629f419972e7ec7b5fdca}.
\end{center}

The verifier establishes the finite polynomial identity exactly.  It does not
reprove the published structural theorems used in
Proposition~\ref{prop:reduction}.  This working proof has not yet undergone
independent peer review.

\begin{thebibliography}{9}
\bibitem{Hwang1986}
S.-G. Hwang,
\emph{A note on a conjecture on permanents},
Linear Algebra and its Applications 76 (1986), 31--44.

\bibitem{CheonWanless2012}
G.-S. Cheon and I. M. Wanless,
\emph{Some results towards the Dittert conjecture on permanents},
Linear Algebra and its Applications 436 (2012), 791--801.
\end{thebibliography}

\end{document}
